\documentclass[11pt,a4paper]{article}
\usepackage{xeCJK}
\usepackage{fontspec}
\usepackage{xcolor}
\usepackage{fancyhdr}
\usepackage{booktabs}
\usepackage{array}
\usepackage{geometry}
\usepackage[protrusion=true]{microtype}
\geometry{margin=1.8cm,top=2cm,headheight=15pt}
\linespread{1.05}

\setmainfont{Baskerville}
\setCJKmainfont{Songti SC}
\newfontfamily{\starfont}{Songti SC}

\definecolor{wrongred}{RGB}{198,40,40}
\definecolor{wrongbg}{RGB}{253,235,233}
\definecolor{rulegray}{RGB}{200,200,200}

\setlength{\emergencystretch}{3em}
\setlength{\fboxsep}{3pt}
\setlength{\parindent}{0pt}

\newcommand{\blank}{\underline{\hspace{2.8cm}}}
\newcommand{\checkmarkbox}{\fbox{\rule{0pt}{0.8ex}\rule{0.8ex}{0pt}}}
\newcommand{\STAR}{\textcolor{wrongred}{\starfont ★}}
% 错答（红色）
\newcommand{\W}[1]{\textcolor{wrongred}{#1}}
% 正确答案（加粗）
\newcommand{\A}[1]{\textbf{#1}}

\pagestyle{fancy}
\fancyhf{}
\fancyhead[L]{\footnotesize 英语单词默写 anyhow -- assure（错词订正）}
\fancyhead[R]{\footnotesize 赵文瀚}
\fancyfoot[C]{\thepage}
\renewcommand{\headrulewidth}{0.4pt}
\renewcommand{\footrulewidth}{0pt}

\begin{document}

\begin{center}
{\LARGE\bfseries 英语单词默写 anyhow -- assure}\\[0.25em]
{\large 错词订正与复盘}\\[0.35em]
{\footnotesize 2029 届高一英语词汇手册默写（pp.21--30，anyhow--assure）· 2026-09}
\end{center}

\vspace{-0.2em}
\noindent{\color{rulegray}\rule{\textwidth}{0.8pt}}

\vspace{0.7em}
\noindent\fcolorbox{black!15}{black!4}{%
  \parbox{\dimexpr\linewidth-2\fboxsep-2\fboxrule\relax}{%
    \small 日期：2026-09\qquad 来源：作业默写（词汇手册 pp.21--30，共 50 词）\qquad 姓名：赵文瀚（班级 6，学号 42）\\[0.3em]
    原卷得分：\textcolor{wrongred}{\textbf{28}}\,/\enspace 50\qquad
    错词：\textcolor{wrongred}{\textbf{22}}\,/\enspace 50\qquad
    重做日期：\underline{\hspace{2.4cm}}}}

\vspace{0.8em}
\noindent 说明：下表按原卷题号列出被判错的 22 个词。\STAR\ 表示该题不仅写错，还暴露出词形/词义的系统问题（拼写、词性、近义混淆），需优先攻破。

\vspace{0.5em}
\renewcommand{\arraystretch}{1.28}
\noindent{\small\begin{tabular}{@{}p{0.75cm}@{\hspace{0.12cm}}p{3.9cm}@{\hspace{0.12cm}}p{4.4cm}@{\hspace{0.12cm}}p{2.7cm}@{\hspace{0.12cm}}p{1.8cm}@{\hspace{0.12cm}}p{2.35cm}@{\hspace{0.12cm}}c@{}}
\toprule
\textbf{题号} & \textbf{中文释义} & \textbf{我的错答} & \textbf{正确答案} & \textbf{错误类型} & \textbf{二刷} & \textbf{掌握} \\
\midrule
6  & 显然的            & \W{appearant}          & \A{apparent}          & 拼写错误 & \checkmarkbox 过\,\checkmarkbox 未过 & 1 \\
8  & 有吸引力 \STAR    & \W{attractive}         & \A{appeal}            & 词义混淆 & \checkmarkbox 过\,\checkmarkbox 未过 & 1 \\
9  & 表面看来 \STAR    & \W{appeal（疑，见注）} & \A{apparently}        & 形近混淆 & \checkmarkbox 过\,\checkmarkbox 未过 & 1 \\
12 & 看来好像 \STAR    & \W{apple／appeal（疑，见注）} & \A{appear}     & 形近混淆 & \checkmarkbox 过\,\checkmarkbox 未过 & 1 \\
13 & 器械 \STAR        & \W{arch（串了 23 题）} & \A{appliance}         & 串行 & \checkmarkbox 过\,\checkmarkbox 未过 & 1 \\
15 & 使（自己）致力于  & \W{（空白）}           & \A{apply}             & 完全没背到 & \checkmarkbox 过\,\checkmarkbox 未过 & 0 \\
19 & 方法              & \W{（空白）}           & \A{approach}          & 完全没背到 & \checkmarkbox 过\,\checkmarkbox 未过 & 0 \\
20 & 近似              & \W{（空白）}           & \A{approximately}     & 完全没背到 & \checkmarkbox 过\,\checkmarkbox 未过 & 0 \\
21 & 批准 \emph{vt.} \STAR & \W{apply（串了 15 题）} & \A{approve}      & 串行 & \checkmarkbox 过\,\checkmarkbox 未过 & 1 \\
22 & 欣赏 \emph{n.}    & \W{（空白）}           & \A{appreciation}      & 完全没背到 & \checkmarkbox 过\,\checkmarkbox 未过 & 0 \\
26 & 算术 \emph{n.}    & \W{arithmati}          & \A{arithmetic}        & 拼写错误 & \checkmarkbox 过\,\checkmarkbox 未过 & 1 \\
28 & 出现 \emph{vi.} \STAR & \W{appear}         & \A{arise}             & 词义混淆 & \checkmarkbox 过\,\checkmarkbox 未过 & 1 \\
29 & 武器（总称）      & \W{arm（漏 s）}        & \A{arms}              & 词形错误 & \checkmarkbox 过\,\checkmarkbox 未过 & 1 \\
30 & 筹备 \emph{n.}    & \W{（空白）}           & \A{arrangement}       & 完全没背到 & \checkmarkbox 过\,\checkmarkbox 未过 & 0 \\
36 & 阿司匹林          & \W{asprin}             & \A{aspirin}           & 拼写错误 & \checkmarkbox 过\,\checkmarkbox 未过 & 1 \\
37 & 指派 \emph{vt.} \STAR & \W{appoint}        & \A{assign}            & 词义混淆 & \checkmarkbox 过\,\checkmarkbox 未过 & 1 \\
38 & 集合 \emph{v.} \STAR & \W{assume}          & \A{assemble}          & 形近混淆 & \checkmarkbox 过\,\checkmarkbox 未过 & 1 \\
43 & 运动的            & \W{（空白）}           & \A{athletic}          & 完全没背到 & \checkmarkbox 过\,\checkmarkbox 未过 & 0 \\
44 & 护理 \emph{v.}    & \W{（空白）}           & \A{attend}            & 完全没背到 & \checkmarkbox 过\,\checkmarkbox 未过 & 0 \\
45 & 达到（某年龄，水平等） & \W{（空白）}      & \A{attain}            & 完全没背到 & \checkmarkbox 过\,\checkmarkbox 未过 & 0 \\
46 & 系，贴            & \W{（空白）}           & \A{attach}            & 完全没背到 & \checkmarkbox 过\,\checkmarkbox 未过 & 0 \\
50 & 确保 \STAR        & \W{attain（串了 45 题）} & \A{assure}        & 串行 & \checkmarkbox 过\,\checkmarkbox 未过 & 1 \\
\bottomrule
\end{tabular}}

\vspace{0.4em}
\noindent{\footnotesize 注：错答栏由原卷扫描件比对所得。第 9、12、13 题连笔较重，经逐字放大比对（字母宽度、下伸笔画数、高字母位置）判定：第 13 题所写宽度 54\,pt 与第 23 题 \emph{arch}（56\,pt）一致、末字母为高字母 \emph{h}，故定为 \emph{arch}（串写第 23 题答案）；第 9 题为 6 字母、末字母为高字母，定为 \emph{appeal}；第 12 题为 5--6 字母、含两个 \emph{p} 下伸笔画，在 \emph{apple}／\emph{appeal} 之间无法确证，请对照 {\xeCJKsetup{CJKecglue={}}\texttt{vocab\_anyhow\_assure\_01\_原卷\_1.jpg}} 复核。第 21 题所写 \emph{apply} 实为第 15 题答案、第 50 题所写 \emph{attain} 实为第 45 题答案，属串行（该两题原卷为空）。}

\newpage
\begin{center}
{\Large\bfseries 错因归类与复盘}
\end{center}
\vspace{-0.4em}
\noindent{\color{rulegray}\rule{\textwidth}{0.8pt}}

\vspace{0.5em}
\noindent 说明：按错误性质归类 22 个错词，供汇总对账；逐词二刷结果见第 1 页错词表「二刷」列。

\vspace{0.5em}
\renewcommand{\arraystretch}{1.35}
\begin{center}
\begin{tabular}{@{}p{4.6cm} p{10.3cm} c@{}}
\toprule
\textbf{错误类型} & \textbf{题号} & \textbf{词数} \\
\midrule
完全没背到（空白） & {\small 15, 19, 20, 22, 30, 43--46} & 9 \\
拼写错误 & {\small 6, 26, 36} & 3 \\
词形错误（词性／单复数） & {\small 29} & 1 \\
词义混淆（写成他词） & {\small 8, 28, 37} & 3 \\
形近混淆（写成形近词） & {\small 9, 12, 38} & 3 \\
串行（写了别题答案） & {\small 13, 21, 50} & 3 \\
\midrule
\textbf{合计} & \textbf{22 题} & \textbf{22} \\
\bottomrule
\end{tabular}
\end{center}

\vspace{0.5em}
\noindent\textbf{关联知识点标签}
\vspace{0.25em}

\noindent\begin{tabular}{@{}p{8.2cm}@{\hspace{0.4cm}}p{8.2cm}@{}}
\begin{minipage}[t]{\linewidth}
\small
\STAR\ \textbf{ap- 词族连锁记忆}：appeal / appear / appearance / apparent / apparently\\[0.2em]
\STAR\ \textbf{形近动词}：appoint（任命）/ assign（指派）；assume（假定）/ assemble（集合）\\[0.2em]
\STAR\ \textbf{拼写易错长词}：apparent（不是 appearant）、arithmetic、aspirin
\end{minipage} &
\begin{minipage}[t]{\linewidth}
\small
\STAR\ \textbf{近义辨析}：appeal（有吸引力）/ attractive（有吸引力的）；arise（出现）/ appear（出现）\\[0.2em]
\STAR\ \textbf{词性意识}：approve（vt. 批准）/ approval（n.）；apply / application / applicant\\[0.2em]
\STAR\ \textbf{单复数}：arms（武器总称）≠ arm（手臂）\qquad
\STAR\ \textbf{范围}：A 字母 ap- / ar- / as- / at- 词族
\end{minipage} \\
\end{tabular}

\vspace{0.65em}
\noindent\fcolorbox{black!15}{black!4}{%
  \parbox{\dimexpr\linewidth-2\fboxsep-2\fboxrule\relax}{%
    \textbf{复盘记录}\\[0.35em]
    最薄弱的板块：\underline{\hspace{9.0cm}}\\[0.55em]
    主要错误原因：\checkmarkbox 完全没背到\quad \checkmarkbox 词形记错\quad \checkmarkbox 拼写不牢\quad \checkmarkbox 词义混淆\quad \checkmarkbox 词组漏词\\[0.55em]
    本次复习的改进措施：\underline{\hspace{10.1cm}}\\[0.55em]
    \hspace{3.2em}\underline{\hspace{10.1cm}}\\[0.55em]
    下次复习日期：\underline{\hspace{2.2cm}}\qquad
    当前掌握度：\checkmarkbox 仍薄弱\quad \checkmarkbox 基本掌握\quad \checkmarkbox 已掌握
  }}

\end{document}
